% AUTORIDAD DE FUENTE: Exposé IX, Proposición 7.5 y §7.6,
% líneas 2082--2103; índices de escaneo 423--424, folios 412--413.
% Las capas francesas D/C no presentan diferencias materiales.

Resumamos, desde el punto de vista de los modelos de Néron, los
resultados obtenidos acerca de las partes abelianas y sus
emparejamientos.

\medskip
\noindent
\underline{Proposición} 7.5. \underline{Sea $S$ un rasgo completo de
cuerpo de funciones $K$. A todo esquema abeliano $A_K$ sobre $K$ con
reducción semiestable sobre $S$, el procedimiento de 7.1, aplicado a la
componente neutra $A^o$ del modelo de Néron $A$ de $A_K$, asocia su
parte abeliana $B=A^{ab}$, que es un esquema abeliano sobre $S$ que
depende funtorialmente de $A_K$. Para cualquier otro esquema abeliano
$A'_K$ con reducción semiestable, se puede asociar a toda correspondencia
divisorial $W_K$ sobre $(A_K,A'_K)$ una única correspondencia divisorial
$W^{ab}$ sobre $(A^{ab},{A'}^{ab})$ tal que, para todo entero
$n\geq0$, la imagen inversa de $W_n^{ab}$ por
$A_n^o\longrightarrow B_n=A_n^{oab}$ y
${A'}_n^o\longrightarrow B'_n={A'}_n^{oab}$ sea la biextensión $W_n$
de $(A_n^o,{A'}_n^o)$ obtenida de $W$ por reducción, donde $W$ es la
extensión canónica (1.4) de $W_K$ a $(A^o,{A'}^o)$.}

%%%%%%%%%%  scan idx 424  |  folio 413  %%%%%%%%%%

\underline{La formación de} $W^{ab}$ \underline{es compatible con la formación de
las imágenes inversas de correspondencias por morfismos de esquemas
abelianos} $\overline A_K\longrightarrow A_K$ \underline{y}
$\overline {A'}_K\longrightarrow A'_K$. \underline{Si} $W_K$ \underline{es una dualidad
perfecta}, \underline{entonces} $W^{ab}$ \underline{también lo es}. \underline{Si} $W_K$ \underline{es la correspondencia
divisorial sobre} $(A_K,A_K)$ \underline{asociada a una polarización de} $A_K$,
\underline{entonces} $W^{ab}$ \underline{está asociada a una polarización de} $A^{ab}$.

La compatibilidad de la formación de los $W^{ab}$ con las imágenes
inversas es inmediata a partir de la caracterización indicada de los
$W^{ab}$. La afirmación adicional relativa a las polarizaciones se
estableció en 7.1.5 al mismo tiempo que la algebraizabilidad de
$\widehat B$.

\medskip
\noindent
7.6. Sigamos suponiendo que $S$ es un rasgo completo, siendo $A$ el
modelo de Néron de un esquema abeliano $A_K$ sobre $K$. Pongamos

\begin{equation}\tag{7.6.1}
A_K^{\natural}=(A^{\natural})_K \quad ,
\end{equation}
donde $A^{\natural}$ es el esquema en grupos de Raynaud (7.2.2): es una
extensión del grupo finito étale $\Phi$ por $A^{\natural o}$, mientras
que $A^{\natural o}$ es una extensión del esquema abeliano $B$ por el
toro $T$. Así, el propio $A_K^{\natural}$ es una extensión de $\Phi_K$
por
$(A_K^{\natural})^o=(A^{\natural o})_K=(A^{o\natural})_K$, que es una
extensión del esquema abeliano $B_K$ por el toro $T_K$. Para todo número
primo $\ell$, (7.3.5) y (7.3.2) dan

\begin{equation}\tag{7.6.2}
T_{\ell}(A_K^{\natural}) \simeq T_{\ell}(A_K^{\natural o})
 \simeq T_{\ell}(A_K)^f \; , \qquad
T_{\ell}(T_K) \simeq T_{\ell}(A_K)^t \quad .
\end{equation}

\medskip
\noindent
7.6.3. Por 3.3, la formación de $A_K^{\natural o}$ conmuta con todo
cambio de base asociado a un morfismo de rasgos completos
$S'\longrightarrow S$.
